% ---
% id: protocol-compiler/def:chaincond
% kind: def
% title: "Interface conditions on a chain"
% summary: >-
%   A presented chain must match descriptors, map physical outcomes onto
%   emitted ones, deliver an exactly admissible target measurement and, where
%   channels are data, satisfy the pullback and source identities exactly. A
%   non-empty source segment adds registry checks and two admission rules.
% prominence: standard
% status: implemented
% uses: [def:declaration, def:validator, def:exact, transmitter/def:source-segment]
% mentions: [thm:chain-consumption, def:gate-outcome-table, def:spec, def:homog, reductions/def:chain, reductions/rem:ordering, reductions/thm:emitted-product, transmitter/def:srcert, reductions/lem:precomposition, reductions/thm:two-sided, rem:source-segment-rules, req:chain-refusals, foundations/req:exact-island, foundations/req:who-snaps, ass:link]
% assumes: []
% refs:
%   - "tupkary2026decoy, Rem. 9.2"
% provenance: newly-derived
% implements: certiqs_key.protocol_compiler.chain:check_chain
% lean:
% notes: "reductions/def:chain is a mention and not a use: the conditions are stated in this module's words on the presented record, and a uses edge would carry the producer's definitional ledger into thm:chain-consumption, thm:band-a and thm:gates. The transmitter edge is a use: (S1)-(S4) are decided here as that definition states them, and it assumes nothing."
% ---

The chain is the compiler's third input. A \emph{squashing} link replaces
Bob's measurement and therefore acts on the receiver space; a
\tref{source link} replaces Alice's signal ensemble and therefore acts on the
\tref{transmitted system} $\Asig$, not on the receiver space at all. Every
link carries two descriptors, of what it starts from and of what it
produces: for a source link these are the transmitter's
\tref{source descriptor} and \tref{target descriptor}; for a receiver link
they are receiver-side, and the last receiver link produces the emitted
system with its \tref{sector map}. One list of links cannot carry both kinds,
and a single descriptor rule cannot type both, so the record is a pair of
segments with a condition set each, and the composition-order rule is a
property of the pair rather than a tag. The warrant for that order is the
published one, \cite[Rem.~9.2]{tupkary2026decoy}, and on the module side
\xref{reductions/rem:ordering}, which states why a source link after a
receiver link is not free; this module types the record and states neither.

That the identity decided at (C4) and (S3) is the identity written is
\xref{foundations/req:exact-island} at work, and a producer repairs its own
emission under \xref{foundations/req:who-snaps}. (C1)--(C5), (S1)--(S4),
the licensed-only rule and the source-mode rule are decided where the chain
is consumed, before a declaration exists, so a failure produces a refusal and
no declaration (Requirement~\ref{req:chain-refusals}), and the \tref{gate outcome
table} of Definition~\ref{def:gate-outcome-table} needs no row for them. The grounds of the
licensed-only rule and of the source-mode rule are
Remark~\ref{rem:source-segment-rules}.

\begin{definition}\label{def:chaincond}
A \tref{reduction chain} $\pcChain=(\Chsrc,\Chrec)$ presented to the compiler, with
$\Chsrc=(\Xi_1,\dots,\Xi_{L_{\mathrm s}})$ and
$\Chrec=(\Lambda_1,\dots,\Lambda_L)$, satisfies the
\term[The conditions a reduction chain presented to the compiler must satisfy, on each of its two segments and between them.]{interface conditions on a chain}
when the following hold.

\emph{\tref[receiver segment]{Receiver segment}.}
\begin{enumerate}[label=(C\arabic*),leftmargin=3.0em]
\item \emph{(descriptor match)} the target descriptor of $\Lambda_i$ equals the
source descriptor of $\Lambda_{i+1}$ for $1\le i<L$; the source descriptor of
$\Lambda_1$ is the physical receiver space of $\Spec_{\mathrm{rec}}$; and the
target descriptor of $\Lambda_L$ is $(\Qsys,\Smap)$ --- that is, the last
link's target descriptor \emph{is} the sector map. That equality is
decided by the same four-component comparison V1 makes on the map of a
declared $\Qsys$-operator family (Definition~\ref{def:validator});
\item \emph{(outcome correspondence)} the declared outcome correspondence is a
tabulated total surjection $\gamma:\Vph\to\Vset$ from the \emph{physical}
alphabet onto the \emph{emitted} one --- the coarse-graining direction, of
which a bijection is the special case;
\item \emph{(\tref[target povm]{target POVM})} $\{F_v\}_{v\in\Vset}$ is a family of
\tref[block list]{block lists} against $\Smap$ --- \emph{typing}, held by the
emitted form and not checked --- and for each $s$ the block family
$\{F_v^{(s)}\}_v$ is a POVM on $\Qsys_s$, decided \emph{exactly} on the
emitted blocks themselves: $F_v^{(s)}$ Hermitian and positive semidefinite
for each $v$, $\sum_vF_v^{(s)}=\idon{\Qsys_s}$, and the declared-order sum
forming it rounding nowhere (Definition~\ref{def:exact}). This condition is
check V3 restricted to $\{F_v\}$ and states the same list of conjuncts
(Theorem~\ref{thm:chain-consumption}\,(ii));
\item \emph{(pullback, where checkable)} if every link of $\Chrec$ supplies its
channel as data then, with
$\Lambda\coloneqq\Lambda_L\circ\dots\circ\Lambda_1$,
\begin{equation}\label{eq:pullback}
  \Lambda^{\dagger}(F_v)=\sum_{w\in\gamma^{-1}(v)}\Fph_w
  \qquad\text{for every }v\in\Vset ,
\end{equation}
an \emph{exact} matrix identity, decided entry by entry on the emitted matrices
and on the supplied channel data, with the fibre sum accumulated in the
\tref{declared order} of $\Vph$. Where some link supplies no channel as data,
only a theorem that constructs one, the identity is not decidable from the
emitted data and is carried instead as an obligation of the chain's
\tref{producer}, Assumption~\ref{ass:link};
\item \emph{(slots)} each link carries its constraint entries and its licence
slot, both possibly empty, and its status tag.
\end{enumerate}

\emph{\tref[source segment]{Source segment}.} The source segment satisfies the
\tref{interface conditions} (S1)--(S4) of
\xref{transmitter/def:source-segment} --- descriptor match with the
\tref{virtual transmitted system} $\Avirt$ as the source descriptor of
$\Xi_1$, on which the certificate of \xref{transmitter/def:srcert} is
written, and every descriptor Alice-side; index preservation; the source
identity $\Xi(\rvirt_x)=\rho_x$ for every $x$, that definition's equation
(S3), where every link supplies its channel as data; and the slots ---
\emph{decided by this module on the presented data}, exactly as that
definition states them. When $L_{\mathrm s}=0$ the segment is empty, $\Xi$ is
the identity, $\Avirt=\Asig$ and $\rvirt_x=\rho_x$, so (S3) holds by
inspection.

\emph{Between the segments.}
\begin{enumerate}[label=(X\arabic*),leftmargin=3.0em]
\item \emph{(order)} the chain is composed source segment first and receiver
segment second. This is not a tag and is not checked: it is the shape of the
record. A source-map link has no slot in $\Chrec$ to be written into, and a
link presented in the wrong segment fails that segment's descriptor rule ---
(C1) for a source map, since $\Asig$ is not the physical receiver space, and
(S1) for a squashing map, since $\Qsys$ is not an Alice-side descriptor.
\end{enumerate}

\emph{At which reading (C4) and (S3) are decided.} Both are \emph{matrix
equalities on \tref{declared data}} --- a composite channel supplied as
binary64 data, applied to binary64 POVM elements on one side and to binary64
signal states on the other --- and both are decided \emph{exactly}, entry by
entry, on the numbers the compiler holds. \emph{The composite, the adjoint
and the fibre sum are formed in exact arithmetic as well}, so the identity
decided is the identity written; that
the operands are binary64 is not what makes the decision exact, a product of
two \tref{binary64 numbers} need not be one. Exactness here is an obligation of
the chain's producer (Assumption~\ref{ass:link}): a producer whose composite
lands off the identity by a rounding has not delivered the reduction the
interface asks for, and repairs its own emission; no window is opened here. Everything else
in this definition --- descriptors, alphabets, fibres, slots, segments --- is
combinatorial and is read exactly, as D(vi)'s relation symbols are.

\emph{The general format: one receiver segment per \tref{measured factor}.}
For a chain presented as Definition~\ref{def:spec}\,(c$'$), with
$\Chrec=(\pcChain^{f}_{\mathrm{rec}})_{f\in\mathbb{F}}$, the conditions
(C1)--(C5) are decided \emph{of each segment separately, and of nothing
across two segments}. (C1): within segment $f$ the target descriptor of a link
is the source descriptor of the next, the source descriptor of
$\Lambda^{f}_1$ is factor $f$'s physical receiver system, and the target
descriptor of $\Lambda^{f}_{L_f}$ is
$(\Qsys_f,\SecSet^{f},(d_{f,s_f})_{s_f})$ --- \emph{that factor's own sector
map, and not the product map}, which is derived by
Definition~\ref{def:declaration}\,(DX1) and is the target descriptor of
nothing; the same four-component comparison decides it. (C2): one tabulated total surjection
$\gamma_f$ from factor $f$'s physical alphabet onto the alphabet factor $f$'s
last link delivers. (C3): the \tref{target family} $\{F^{f}_{v_f}\}$ is a family of
block lists against factor $f$'s map and, for each $s_f$, a POVM on
$\Qsys_{f,s_f}$ decided exactly on that factor's blocks --- which is check V3
restricted to that factor, conjunct for conjunct, in its general-format
reading. (C4): where every link of segment $f$ supplies its channel as data,
with $\Lambda_f\coloneqq\Lambda^{f}_{L_f}\circ\dots\circ\Lambda^{f}_1$,
\begin{equation}\label{eq:pullback-factor}
  \Lambda_f^{\dagger}\big(F^{f}_{v_f}\big)
  =\sum_{w_f\in\gamma_f^{-1}(v_f)}F^{\mathrm{ph},f}_{w_f}
  \qquad\text{for every }v_f ,
\end{equation}
an exact matrix identity on that factor's matrices, the fibre sum accumulated
in the declared order of that factor's physical alphabet. \emph{The product
composite $\bigotimes_f\Lambda_f$ and the product identity on $\{F_v\}$ are
formed and decided by no one}: the identity
$\big(\bigotimes_f\Lambda_f\big)^{\dagger}(F_v)
=\bigotimes_f\Lambda_f^{\dagger}(F^{f}_{v_f})$ is a consequence of
\eqref{eq:pullback-factor} and of the product rule
\eqref{eq:productfamily} --- it is \xref{reductions/thm:emitted-product}\,(ii),
named here and decided here by nothing --- and (C4) is the per-factor
identity and nothing wider. (C5) is per link, as above. Where some link of
segment $f$ supplies no channel as data, only a theorem that constructs one,
\eqref{eq:pullback-factor}
for \emph{that} segment is not decidable from the emitted data and is carried
as Assumption~\ref{ass:link}, exactly as above and per factor.

\emph{A non-empty source segment, and the conditions it brings.} \emph{A
presented chain may carry one.} The rule by which a non-empty source segment
composes with a receiver segment is a theorem of the reductions module
(\xref{reductions/lem:precomposition}, \xref{reductions/thm:two-sided}, with
\xref{reductions/def:chain}), in
\tref{trusted-source mode}; and the module that admits source-link kinds may
emit an artefact carrying a link of an admitted kind. What follows is what
this module decides of such a segment, on the presented data and exactly;
whether a presented chain carries one is data of the record and is asserted
here of no artefact.
\begin{itemize}[leftmargin=2.6em]
\item \emph{(S1)--(S4) as \xref{transmitter/def:source-segment}
states them}, which is what the source-segment paragraph above already says;
and within (S3), \emph{the split by kind of witness}: for a link supplying its
channel as data the per-link identity is decided here, entry by entry, with
the composite formed in exact arithmetic; for a link supplying a
\tref{theorem constructor} it is \emph{not decided here at all} but is the conclusion of the
claim that link's \tref{obligation record} names, and the composite identity
is the composition of the per-link ones. So what this module decides of (S3)
is the data links' share of it, and the constructor links' share travels as
the obligation Assumption~\ref{ass:link} names ---
Theorem~\ref{thm:chain-consumption}\,(iv), (v) --- and is discharged nowhere in this
contract.
\item \emph{(A1)--(A3) of that definition, which are decided on data and are
not part of (S1)--(S4)}: that the triple a link's obligation record carries
--- its kind, the claim id it names, and the assumption short ids it names ---
appears \emph{entry by entry} in the \tref{soundness registry} the artefact
carrying the segment emits, both components of that registry entry's
assumption list compared, so that a record naming a claim's closure and
omitting the \tref{source model}'s attestation fails as one naming neither
does; that every assumption of that entry is in the artefact's assumption
ledger; and that the segment is presented inside an artefact of the module
that admits its links. These are comparisons of finitely many declared
strings and of a finite list against a finite list, and \emph{they are decided
here because they are decidable on the presented record}, which is why they
are named beside (S1)--(S4) rather than left to Assumption~\ref{ass:link}.
\emph{What they do not decide} is that the named claim's conclusion applies to
these descriptors --- family membership is the emitting module's assertion and
no consumer re-proves a theorem --- and a link supplying its channel as data
carries no obligation record and is decided by (S3) instead.
\item \emph{the licensed-only rule.} Beside a non-empty source segment
\emph{every} entry of the emitted $\CList$ is required to carry a real scalar
as its right-hand side --- homogeneous, party-local, or an untrusted-mode
scalar (Definition~\ref{def:homog}).
\item \emph{the declared source mode.} \emph{A declaration carrying a
non-empty source segment is in trusted-source mode}: the rule fails
where the \tref{marginal slot} of \eqref{eq:declaration} holds $\nomarg$
(Definition~\ref{def:declaration}), the \tref{untrusted-source mode}, while
$L_{\mathrm s}\ge1$.
\end{itemize}
\emph{(X1) in the general format.} The order rule is unchanged and is still
the type of the record: a source-map link has no slot in any receiver segment
to be written into, and a link presented in the wrong list fails that list's
descriptor rule --- (C1) of the factor whose segment it was written into, or
(S1). In the reductions module the composite of a chain is a pair, and the
reason a source link may not follow a receiver link is a theorem of that
module (\xref{reductions/lem:precomposition},
\xref{reductions/thm:two-sided}) beside the published warrant
\cite[Rem.~9.2]{tupkary2026decoy}; this module types the record, decides the
descriptors, and states neither.
\end{definition}
